% !TeX root = ../main.tex
\section{Related Work}
\label{sec:related_work}

\subsection{Harmful Meme Detection and Structured Understanding}
\label{sec:related:detection}

Harmful and hateful meme detection is inherently multimodal because the
harmful meaning may emerge only from the interaction between an image and
its accompanying text.
The Hateful Memes benchmark~\cite{kiela2020hateful} helped establish this setting, and subsequent
work has explored stronger multimodal fusion and pretrained
vision--language representations.
CLIP-based methods such as HateCLIPper~\cite{kumar2022hateclipper} and
MemeCLIP~\cite{shah2024memeclip} model image--text
interactions for meme classification, while MOMENTA combines global and
local multimodal evidence and extends detection toward harmful-target
identification~\cite{pramanick2021momenta}.
Related approaches such as DISARM~\cite{sharma2022disarm} likewise emphasize identifying victims
or targets rather than predicting harmfulness alone.

Recent work increasingly moves beyond a single binary label.
SAFE-MEME~\cite{nandi2024safememe} uses structured question--answer reasoning and hierarchical
categorization, including questions about the target and type of hate.
EXPLAINHM~\cite{lin2024explainhm} and
EXPLAINHM++~\cite{lin2026explainhmpp} generate natural-language rationales through
multimodal debate, and recent large-multimodal-model approaches such as
RA-HMD~\cite{mei2025rahmd} and ExPO-HM~\cite{mei2026expohm} incorporate fine-grained supervision, reasoning, or
explanation into harmful-meme modeling.
These developments motivate richer outputs, but target prediction,
fine-grained categorization, rationale generation, and evidence grounding
are still often represented as method-specific objectives.
Our formulation instead places harmfulness, target presence and
granularity, primary intent, rhetorical tactics, and multimodal relation in
one structured prediction space, while retaining a separate evidence record
for audit.

\subsection{Knowledge-Augmented Meme Reasoning}
\label{sec:related:knowledge}

Meme interpretation can depend on social, cultural, and event-specific
knowledge that is absent from the pixels or OCR text.
PromptHate introduces image-derived textual context into prompt-based
classification, whereas retrieval-based methods explicitly introduce
related examples or documents~\cite{cao2022prompthate}.
RGCL~\cite{mei2024rgcl} uses retrieved meme relationships for contrastive representation
learning and retrieval-based inference, and RA-HMD combines multimodal
adaptation with retrieval-oriented representations for out-of-domain
classification.
EXPLAINHM++ retrieves related harmful and harmless memes to ground a second
round of multimodal debate before producing its final decision and
rationale.

A closely related issue is robustness to noisy retrieval.
MerFT~\cite{lee2026merft} studies distractor-aware multimodal meme
reasoning and trains retrieval-augmented models to remain robust to
semantically plausible but misleading documents.
Its experiments explicitly vary distractor frequency, demonstrating that
retrieval quality can materially affect meme reasoning.
This reinforces the need to treat retrieved context as potentially noisy
rather than inherently trustworthy.

Our focus is complementary but different.
Instead of relying only on retrieval-aware training to become robust to
distractors, we place an explicit candidate-level verifier between
knowledge acquisition and multimodal fusion.
Each external candidate is evaluated against internal meme evidence for
relevance, validity, and task-specific support before it can influence the
structured prediction.
The framework additionally retains both accepted and rejected provenance,
allowing retrieval quality and verification behavior to be analyzed
separately.
Thus, the central distinction of our approach is not retrieval itself, but
the explicit conversion of acquired context into \emph{verified evidence}
before task-aware reasoning.
